\begin{recipe}{Frituursnacks uit de diepvries}
\kicker[Bijgerecht,Nederlands,Snack]{Bijgerecht · Nederlands · snack}
\index[register]{Frituursnacks uit de diepvries}
\index[register]{Kroket}
\index[register]{Frikandel}
\index[register]{Bitterbal}
\index[register]{Kaassouffl\'e}
\index[register]{Friet}

\meta{3 tot 8 minuten \dvd Frituurvet 175--180\,°C}

\lettrine{G}{een} echt recept, meer een geheugensteun: wij hebben geen
frituurpan in huis, maar staan bij toerbeurt soms wel een middag te
frituren bij de kantine van de honkbal. Dit zijn de tijden voor de
standaard diepvriessnacks, direct uit de diepvries de olie in.

\blockrule{Ingrediënten}
\begin{ingredients}
  \ingb{frituurolie of frituurvet, ruim}
  \ingb{kroketten}
  \ingb{frikandellen}
  \ingb{bitterballen}
  \ingb{kaassoufflés}
  \ingb{frieten (diepvries, voorgebakken)}
\end{ingredients}

\blockrule{Bereiding}
\tip{Test de temperatuur zonder thermometer met een blokje brood: bruint
het in ongeveer 15 seconden, dan zit de olie rond de 175\,°C.}
\begin{steps}
  \item Verhit de frituurolie tot 175 à 180\,°C. Frituur alles direct
        vanuit de diepvries, niet eerst laten ontdooien.
  \item Kroketten: 4 tot 5 minuten, tot ze goudbruin zijn en vanbinnen
        overal warm.
  \item Frikandellen: 5 tot 6 minuten, tot de korst goudbruin en krokant
        is.
  \item Bitterballen: 3 tot 4 minuten, tot ze goudbruin zijn: ze zijn
        klein en garen sneller dan kroketten.
  \item Kaassoufflés: 4 tot 5 minuten, tot de korst goudbruin is. Laat ze
        even afkoelen voor je erin bijt, de gesmolten kaas vanbinnen is
        heet en kan spatten als de korst lekt.
  \item Frieten: 3 tot 5 minuten, dun gesneden friet is eerder klaar dan
        een dikke steakhouse-snit. Laat alles uitlekken op keukenpapier
        en bestrooi de friet pas met zout zodra hij uit de olie komt.
\end{steps}
\end{recipe}
